% Erratum E15 for inclusion in main.tex.
% Suggested use: % Erratum E15 for inclusion in main.tex.
% Suggested use: % Erratum E15 for inclusion in main.tex.
% Suggested use: % Erratum E15 for inclusion in main.tex.
% Suggested use: \input{erratum_E15_discrepancy_inequality}

\subsection*{Erratum E15: the sufficient inequality in the discrepancy-principle paragraph}
\label{s:erratum:E15}

In the paragraph following Proposition~3.4, the displayed
sufficient condition for the discrepancy principle has the inequality in
the wrong direction.  The corrected condition is
\begin{align}
\label{e:erratum:E15:corrected-condition}
        c_2\,\varepsilon_k^{\nu+1/2}
        \normiii{x^\dag-x_0}{\nu}
        \leq
        \delta(\tau-1).
\end{align}
The printed condition has the two sides of
\eqref{e:erratum:E15:corrected-condition} reversed.

Indeed, using the upper estimate in Proposition~3.4 with
$r=1/2$ gives
\begin{align}
\label{e:erratum:E15:residual-upper-bound}
        \norm{y^\dag-Tx_k}{}
        & =
        \norm{T(x^\dag-x_k)}{}
        =
        \norm{(T^*T)^{1/2}(x^\dag-x_k)}{}
        \\
        & \leq
        c_2\,\varepsilon_k^{\nu+1/2}
        \normiii{x^\dag-x_0}{\nu} .
\end{align}
Therefore \eqref{e:erratum:E15:corrected-condition} implies
\[
        \norm{y^\dag-Tx_k}{}
        \leq
        \delta(\tau-1).
\]
Together with the residual comparison established just before this
paragraph,
\[
        \left|
        \norm{y^\dag-Tx_k}{}
        -
        \norm{y^\delta-Tx_k^\delta}{}
        \right|
        \leq
        \delta,
\]
this yields
\[
        \norm{y^\delta-Tx_k^\delta}{}
        \leq
        \norm{y^\dag-Tx_k}{}+
        \delta
        \leq
        \delta\tau,
\]
which is precisely the discrepancy principle in equation (68).  The printed
reverse inequality does not imply the desired residual estimate; it only
says that the available upper bound in
\eqref{e:erratum:E15:residual-upper-bound} is at least the target
threshold.

The corrected inequality also gives the same iteration-count estimate as
claimed in the text.  For $r=1/2$,
\[
        \varepsilon_k
        =
        \frac{1}{\sigma}
        \frac{\nu+1/2}{k+\nu+1/2} .
\]
Put
\[
        a:=\nu+\frac12,
        \qquad
        R:=\normiii{x^\dag-x_0}{\nu} .
\]
If $R=0$ the assertion is trivial.  Otherwise,
\eqref{e:erratum:E15:corrected-condition} is equivalent to
\[
        c_2 R
        \left(
                \frac{a}{\sigma(k+a)}
        \right)^a
        \leq
        \delta(\tau-1),
\]
or
\[
        k+a
        \geq
        \frac{a}{\sigma}
        \left(
                \frac{c_2 R}{\delta(\tau-1)}
        \right)^{1/a} .
\]
Thus the discrepancy principle is met after at most
\[
        k^\star
        =
        O\!\left(\delta^{-1/a}\right)
        =
        O\!\left(\delta^{-2/(2\nu+1)}\right)
\]
iterations, with constants depending on the fixed quantities
$\sigma$, $\tau$, $\nu$, $c_2$, and
$\normiii{x^\dag-x_0}{\nu}$.  Thus only the direction of the displayed
sufficient inequality has to be corrected.


\subsection*{Erratum E15: the sufficient inequality in the discrepancy-principle paragraph}
\label{s:erratum:E15}

In the paragraph following Proposition~3.4, the displayed
sufficient condition for the discrepancy principle has the inequality in
the wrong direction.  The corrected condition is
\begin{align}
\label{e:erratum:E15:corrected-condition}
        c_2\,\varepsilon_k^{\nu+1/2}
        \normiii{x^\dag-x_0}{\nu}
        \leq
        \delta(\tau-1).
\end{align}
The printed condition has the two sides of
\eqref{e:erratum:E15:corrected-condition} reversed.

Indeed, using the upper estimate in Proposition~3.4 with
$r=1/2$ gives
\begin{align}
\label{e:erratum:E15:residual-upper-bound}
        \norm{y^\dag-Tx_k}{}
        & =
        \norm{T(x^\dag-x_k)}{}
        =
        \norm{(T^*T)^{1/2}(x^\dag-x_k)}{}
        \\
        & \leq
        c_2\,\varepsilon_k^{\nu+1/2}
        \normiii{x^\dag-x_0}{\nu} .
\end{align}
Therefore \eqref{e:erratum:E15:corrected-condition} implies
\[
        \norm{y^\dag-Tx_k}{}
        \leq
        \delta(\tau-1).
\]
Together with the residual comparison established just before this
paragraph,
\[
        \left|
        \norm{y^\dag-Tx_k}{}
        -
        \norm{y^\delta-Tx_k^\delta}{}
        \right|
        \leq
        \delta,
\]
this yields
\[
        \norm{y^\delta-Tx_k^\delta}{}
        \leq
        \norm{y^\dag-Tx_k}{}+
        \delta
        \leq
        \delta\tau,
\]
which is precisely the discrepancy principle in equation (68).  The printed
reverse inequality does not imply the desired residual estimate; it only
says that the available upper bound in
\eqref{e:erratum:E15:residual-upper-bound} is at least the target
threshold.

The corrected inequality also gives the same iteration-count estimate as
claimed in the text.  For $r=1/2$,
\[
        \varepsilon_k
        =
        \frac{1}{\sigma}
        \frac{\nu+1/2}{k+\nu+1/2} .
\]
Put
\[
        a:=\nu+\frac12,
        \qquad
        R:=\normiii{x^\dag-x_0}{\nu} .
\]
If $R=0$ the assertion is trivial.  Otherwise,
\eqref{e:erratum:E15:corrected-condition} is equivalent to
\[
        c_2 R
        \left(
                \frac{a}{\sigma(k+a)}
        \right)^a
        \leq
        \delta(\tau-1),
\]
or
\[
        k+a
        \geq
        \frac{a}{\sigma}
        \left(
                \frac{c_2 R}{\delta(\tau-1)}
        \right)^{1/a} .
\]
Thus the discrepancy principle is met after at most
\[
        k^\star
        =
        O\!\left(\delta^{-1/a}\right)
        =
        O\!\left(\delta^{-2/(2\nu+1)}\right)
\]
iterations, with constants depending on the fixed quantities
$\sigma$, $\tau$, $\nu$, $c_2$, and
$\normiii{x^\dag-x_0}{\nu}$.  Thus only the direction of the displayed
sufficient inequality has to be corrected.


\subsection*{Erratum E15: the sufficient inequality in the discrepancy-principle paragraph}
\label{s:erratum:E15}

In the paragraph following Proposition~3.4, the displayed
sufficient condition for the discrepancy principle has the inequality in
the wrong direction.  The corrected condition is
\begin{align}
\label{e:erratum:E15:corrected-condition}
        c_2\,\varepsilon_k^{\nu+1/2}
        \normiii{x^\dag-x_0}{\nu}
        \leq
        \delta(\tau-1).
\end{align}
The printed condition has the two sides of
\eqref{e:erratum:E15:corrected-condition} reversed.

Indeed, using the upper estimate in Proposition~3.4 with
$r=1/2$ gives
\begin{align}
\label{e:erratum:E15:residual-upper-bound}
        \norm{y^\dag-Tx_k}{}
        & =
        \norm{T(x^\dag-x_k)}{}
        =
        \norm{(T^*T)^{1/2}(x^\dag-x_k)}{}
        \\
        & \leq
        c_2\,\varepsilon_k^{\nu+1/2}
        \normiii{x^\dag-x_0}{\nu} .
\end{align}
Therefore \eqref{e:erratum:E15:corrected-condition} implies
\[
        \norm{y^\dag-Tx_k}{}
        \leq
        \delta(\tau-1).
\]
Together with the residual comparison established just before this
paragraph,
\[
        \left|
        \norm{y^\dag-Tx_k}{}
        -
        \norm{y^\delta-Tx_k^\delta}{}
        \right|
        \leq
        \delta,
\]
this yields
\[
        \norm{y^\delta-Tx_k^\delta}{}
        \leq
        \norm{y^\dag-Tx_k}{}+
        \delta
        \leq
        \delta\tau,
\]
which is precisely the discrepancy principle in equation (68).  The printed
reverse inequality does not imply the desired residual estimate; it only
says that the available upper bound in
\eqref{e:erratum:E15:residual-upper-bound} is at least the target
threshold.

The corrected inequality also gives the same iteration-count estimate as
claimed in the text.  For $r=1/2$,
\[
        \varepsilon_k
        =
        \frac{1}{\sigma}
        \frac{\nu+1/2}{k+\nu+1/2} .
\]
Put
\[
        a:=\nu+\frac12,
        \qquad
        R:=\normiii{x^\dag-x_0}{\nu} .
\]
If $R=0$ the assertion is trivial.  Otherwise,
\eqref{e:erratum:E15:corrected-condition} is equivalent to
\[
        c_2 R
        \left(
                \frac{a}{\sigma(k+a)}
        \right)^a
        \leq
        \delta(\tau-1),
\]
or
\[
        k+a
        \geq
        \frac{a}{\sigma}
        \left(
                \frac{c_2 R}{\delta(\tau-1)}
        \right)^{1/a} .
\]
Thus the discrepancy principle is met after at most
\[
        k^\star
        =
        O\!\left(\delta^{-1/a}\right)
        =
        O\!\left(\delta^{-2/(2\nu+1)}\right)
\]
iterations, with constants depending on the fixed quantities
$\sigma$, $\tau$, $\nu$, $c_2$, and
$\normiii{x^\dag-x_0}{\nu}$.  Thus only the direction of the displayed
sufficient inequality has to be corrected.


\subsection*{Erratum E15: the sufficient inequality in the discrepancy-principle paragraph}
\label{s:erratum:E15}

In the paragraph following Proposition~3.4, the displayed
sufficient condition for the discrepancy principle has the inequality in
the wrong direction.  The corrected condition is
\begin{align}
\label{e:erratum:E15:corrected-condition}
        c_2\,\varepsilon_k^{\nu+1/2}
        \normiii{x^\dag-x_0}{\nu}
        \leq
        \delta(\tau-1).
\end{align}
The printed condition has the two sides of
\eqref{e:erratum:E15:corrected-condition} reversed.

Indeed, using the upper estimate in Proposition~3.4 with
$r=1/2$ gives
\begin{align}
\label{e:erratum:E15:residual-upper-bound}
        \norm{y^\dag-Tx_k}{}
        & =
        \norm{T(x^\dag-x_k)}{}
        =
        \norm{(T^*T)^{1/2}(x^\dag-x_k)}{}
        \\
        & \leq
        c_2\,\varepsilon_k^{\nu+1/2}
        \normiii{x^\dag-x_0}{\nu} .
\end{align}
Therefore \eqref{e:erratum:E15:corrected-condition} implies
\[
        \norm{y^\dag-Tx_k}{}
        \leq
        \delta(\tau-1).
\]
Together with the residual comparison established just before this
paragraph,
\[
        \left|
        \norm{y^\dag-Tx_k}{}
        -
        \norm{y^\delta-Tx_k^\delta}{}
        \right|
        \leq
        \delta,
\]
this yields
\[
        \norm{y^\delta-Tx_k^\delta}{}
        \leq
        \norm{y^\dag-Tx_k}{}+
        \delta
        \leq
        \delta\tau,
\]
which is precisely the discrepancy principle in equation (68).  The printed
reverse inequality does not imply the desired residual estimate; it only
says that the available upper bound in
\eqref{e:erratum:E15:residual-upper-bound} is at least the target
threshold.

The corrected inequality also gives the same iteration-count estimate as
claimed in the text.  For $r=1/2$,
\[
        \varepsilon_k
        =
        \frac{1}{\sigma}
        \frac{\nu+1/2}{k+\nu+1/2} .
\]
Put
\[
        a:=\nu+\frac12,
        \qquad
        R:=\normiii{x^\dag-x_0}{\nu} .
\]
If $R=0$ the assertion is trivial.  Otherwise,
\eqref{e:erratum:E15:corrected-condition} is equivalent to
\[
        c_2 R
        \left(
                \frac{a}{\sigma(k+a)}
        \right)^a
        \leq
        \delta(\tau-1),
\]
or
\[
        k+a
        \geq
        \frac{a}{\sigma}
        \left(
                \frac{c_2 R}{\delta(\tau-1)}
        \right)^{1/a} .
\]
Thus the discrepancy principle is met after at most
\[
        k^\star
        =
        O\!\left(\delta^{-1/a}\right)
        =
        O\!\left(\delta^{-2/(2\nu+1)}\right)
\]
iterations, with constants depending on the fixed quantities
$\sigma$, $\tau$, $\nu$, $c_2$, and
$\normiii{x^\dag-x_0}{\nu}$.  Thus only the direction of the displayed
sufficient inequality has to be corrected.
